\documentclass[10pt, aspectratio=169]{beamer}
% \documentclass[10pt, handout]{beamer}


\usepackage{clases}

\title{Investigación Operativa - Clase 3}
\subtitle{Disyunción}
\date{}
\author{Nazareno Faillace Mullen}
\institute{Departamento de Matemática - FCEN - UBA}

%% Si es versión completa o para la clase
\newif\ifCompleta

% Descomentar para clase completa
\Completatrue

\begin{document}

\maketitle

% \begin{frame}{Contenidos}
%   \setbeamertemplate{section in toc}[sections numbered]
%   \tableofcontents[hideallsubsections]
% \end{frame}

% \section{Introducción}

\begin{frame}{Para ir pensando}
    \large
    Sea $x$ una variable continua o entera (no necesariamente $\geq 0$) y
    $C$ una constante positiva, ¿cómo modelarías linealmente cada una de las siguientes
    restricciones?
    \begin{enumerate}
        \item $|x|\leq C$
        \item $|x|\geq C$
    \end{enumerate}

    \centering
    \textbf{Puesta en común a las 19:30}
\end{frame}

\begin{frame}{Módulo}
    \begin{itemize}
        \item
        Puede ocurrir que agunas restricciones involucren al módulo de algunas variables.
        $$|x|\leq C$$
        donde $C> 0$. Debemos expresarla de manera lineal.

        Modelamos esa expresión con las siguientes restricciones:
        $$\begin{array}{l}
            x\leq C  \\
            x\geq -C
        \end{array}$$
    \end{itemize}
\end{frame}

\begin{frame}{Módulo}
    \begin{itemize}
    \item Lo mismo sirve si tenemos una restricción con módulo en general. Sea $g$ lineal:
    $$|g(x_1,\dots,x_n)| \leq C$$
    Se puede modelar con las restricciones:
    $$\begin{array}{rl}
        g(x_1,\dots,x_n) &\leq C  \\
        -g(x_1,\dots,x_n) &\leq C
    \end{array}$$

    \item O todavía más en general, si $g$ y $h$ son lineales:
    $$|g(x_1,\dots,x_n)| + h(x_1,\dots,x_n)  \leq C$$
    Se puede modelar con:
    $$\begin{array}{rl}
        g(x_1,\dots,x_n) + h(x_1,\dots,x_n)  &\leq C  \\
        -g(x_1,\dots,x_n) + h(x_1,\dots,x_n) &\leq C
    \end{array}$$
    \end{itemize}
\end{frame}

\begin{frame}{Módulo}
    ¿Qué ocurre si la restricción es $|x| \geq c$?
    \pause
    \begin{itemize}
        \item Recordar que:
        \[ |x|\geq c \quad \Leftrightarrow \quad x\geq c \,\lor\, x\leq -c
        \]
    \end{itemize}
\end{frame}

\begin{frame}
    Podemos añadir una variable binaria $\theta$, una constante $M$ lo suficientemente grande y las siguientes
    restricciones:
        \[
        \begin{array}{rl}
            x & \geq c - \theta M  \\
            x & \leq -c + (1-\theta)M \\
            \multicolumn{2}{c}{\theta\in\{0,1\}}
        \end{array}
        \]
\end{frame}


\begin{frame}{Módulo}
    \begin{itemize}
        \item En general, para modelar una restricción del tipo:
        $$|g(x_1,\dots,x_n)| \geq C $$
        Consideramos una variable binaria $\theta$ que nos sirva de indicadora, $M$ un número lo suficientemente grande y a\~nadimos las siguientes restricciones:
        $$\begin{array}{rl}
            g(x_1,\dots,x_n) &\geq C  - \theta M  \\
            g(x_1,\dots,x_n)  &\leq -C + (1-\theta)M \\
            \multicolumn{2}{c}{\theta\in\{0,1\}}
        \end{array}
        $$
    \end{itemize}
\end{frame}

\begin{frame}{Disyunción}
    \begin{itemize}
        \item En general, cuando queremos modelar una disyunción del tipo:
        \[ g_1(x_1,\dots, x_n) \leq b_1 \; \vee \; g_2(x_1,\dots, x_n) \geq b_2
        \]
        con $g_1$ y $g_2$ funciones lineales. Tomando $M$ lo suficientemente grande como para que las condiciones de
        $\leq$ queden satisfechas automáticamente si sumamos $M$ al lado derecho y que las condiciones de $\geq$
        queden automáticamente satisfechas si restamos M del lado derecho, se modela la disyunción de la siguiente
        manera:
        \[\begin{array}{rl}
            g_1(x_1,\dots, x_n) \leq & b_1 + M(1-\theta) \\
            g_2(x_1,\dots, x_n) \geq & b_2 - M\theta \\
            \multicolumn{2}{c}{\theta\in\{0,1\}}
        \end{array}
        \]
        % \textbf{Obs:} $p\Rightarrow q \equiv \lnot p\lor q$
    \end{itemize}
\end{frame}

\begin{frame}{Ejercicio}
    Una planta de producción de MicroApple cuenta con $N$ operarios para realizar $T$ tareas
    distintas. A cada tarea $t$ deben ser asignados una cantidad de operarios en el intervalo
    $[\ell_t, u_t]$ y ningún operario puede quedar sin ser asignado a una tarea.
    Además, debemos tener en cuenta los siguientes requerimientos propios de la producción:
    \begin{enumerate}
        \item hay un conjunto de tareas $\Gamma\subseteq T$, tales que para cada $t\in \Gamma$, el valor absoluto de la
        diferencia entre la cantidad de operarios asignados a las tareas $t$ y $t+1$ debe ser menor o igual que
        un número conocido $m_t$.
        \item debe ocurrir alguna de las siguientes situaciones:
            \begin{enumerate}
                \item la cantidad de operarios asignados a la tarea 1 debe ser al menos el 25\% de los
                asignados a la tarea 2.
                \item se asignan a lo sumo $E$ empleados al conjunto de tareas $\{3,4,5\}$.
            \end{enumerate}
    \end{enumerate}

    Escribir un modelo que permita hallar cuántos operarios pueden asignarse a cada tarea. ¿Cuál es la función
    objetivo?
\end{frame}

\ifCompleta

\begin{frame}
    \textbf{Variables:}

    \begin{tabularx}{\textwidth}{lX}
        $x_t\colon$ & cantidad de operarios asignados a la tarea $t$ ($x_t\in\mathbb{Z}_{\geq 0}$) \\
        $\theta\colon$ & variable auxiliar para la disyunción ($\theta\in\{0,1\}$)
    \end{tabularx}
\end{frame}

\begin{frame}
    \[
        \begin{array}{lrlr}
        \max & 0 &  & \\
        \text{s.a.:} & \displaystyle\sum_{i=1}^{|T|} x_t =& N \\
          & x_t \leq& u_t & t\in T\\
          & x_t \geq & \ell_t & t\in T\\
          & x_t - x_{t+1} \leq & m_t & t\in\Gamma \\
          & x_{t+1} - x_t \leq & m_t & t\in\Gamma \\
          & x_1 \geq & 0.25x_2 - M\theta \\
          & x_3+x_4+x_5 \leq & E + M(1-\theta)\\
          & x_t\in &\mathbb{Z_{\geq 0}} & t\in T \\
          & \theta\in&\{0,1\} \\
        \end{array}
    \]
\end{frame}
\fi

\end{document}



